\subsection{Nasi Padang dan Bentuk Penyajiannya}

Nasi Padang merupakan penyajian nasi putih dengan berbagai pilihan lauk yang berasal dari tradisi kuliner Minangkabau di Sumatra Barat. Pilihan lauk pada rumah makan Padang dapat berupa olahan daging, ayam, ikan, telur, sayuran, dan sambal. Beberapa lauk yang umum dikenal antara lain rendang, ayam goreng, telur, sayur nangka, dan daun singkong \cite{indonesiaTravelPadang,michelinNasiPadang}.

Pada pelayanan rumah makan Padang, terdapat cara penyajian \textit{hidang} dan pelayanan berdasarkan pesanan pelanggan. Pada cara \textit{hidang}, berbagai lauk diletakkan dalam piring-piring kecil di meja pelanggan. Pelanggan memilih lauk yang akan dikonsumsi dan pembayaran ditentukan berdasarkan lauk yang diambil. Pada pelayanan berdasarkan pesanan, pelanggan memilih lauk dari tempat penyajian, kemudian nasi dan lauk yang dipilih disusun pada satu piring atau ditempatkan dalam bungkus untuk dibawa pulang \cite{indonesiaTravelPadang,michelinNasiPadang}.

Bentuk penyajian tersebut perlu dijelaskan karena penelitian ini tidak mencakup seluruh mekanisme pelayanan rumah makan Padang. Sistem yang dirancang hanya ditujukan untuk pelanggan \textit{dine-in} yang mengambil lauk secara \textit{self-service}. Setelah lauk dipilih, seluruh makanan ditempatkan pada satu piring terbuka dan piring tersebut diletakkan pada area pengambilan citra. Sistem tidak digunakan untuk mengenali seluruh piring pada penyajian \textit{hidang} dan tidak mencakup makanan \textit{take-away} yang ditempatkan dalam bungkus atau kotak.

Objek yang dikenali dibatasi pada sembilan kelas, yaitu nasi, rendang, ayam goreng, perkedel, sayur nangka, daun singkong, telur dadar, telur kari, dan sambal. Kesembilan kelas tersebut merupakan batas pengenalan sistem dan tidak dimaksudkan untuk mewakili seluruh variasi lauk yang tersedia pada rumah makan Padang.

